\documentclass[acmsmall,nonacm,screen]{acmart}

\usepackage{mathpartir}
\usepackage{booktabs}

\citestyle{acmauthoryear}
\settopmatter{printacmref=false}

\AtEndPreamble{\theoremstyle{acmdefinition}\newtheorem{remark}[theorem]{Remark}}

\newcommand{\kw}[1]{\mathsf{#1}}
\newcommand{\Str}[2]{\kw{Str}[#1 \triangleright #2]}
\newcommand{\cat}{\mathbin{+\!\!+}}
\newcommand{\recv}[1]{\kw{recv}\,#1}
\newcommand{\hear}[1]{\kw{hear}\,#1}
\newcommand{\tail}[2]{\kw{tail}_{#1}\,#2}
\newcommand{\call}[3]{\kw{call}(#1;\,#2,\,#3)}
\newcommand{\letin}[2]{\kw{let}\;#1 = #2\;\kw{in}}
\newcommand{\letx}[2]{\kw{let}\;#1 = #2}
\newcommand{\iter}[4]{\kw{iter}\;#1\;\kw{from}\;#2\;\kw{as}\;#3.\;#4}
\newcommand{\orr}{\mathbin{\kw{or}}}
\newcommand{\ev}[5]{\langle #1,\, #2\rangle \Downarrow^{#3} \langle #4,\, #5\rangle}
\newcommand{\sem}[1]{[\![#1]\!]}
\newcommand{\fl}[1]{\lfloor #1 \rfloor}
\newcommand{\hit}{\mathrm{hit}}
\newcommand{\pre}{\sqsubseteq}
\newcommand{\eps}{\varepsilon}
\newcommand{\nm}[1]{\mathit{#1}}
\newcommand{\turn}{\nm{turn}}
\newcommand{\hturn}{\nm{hturn}}
\newcommand{\defeq}{\;\triangleq\;}
\newcommand{\tsum}{\textstyle\sum}
\newcommand{\N}{\mathbb{N}}
\newsavebox{\progbox}
\newenvironment{prog}
  {\begin{lrbox}{\progbox}$\displaystyle\begin{array}{@{}l@{\qquad}l@{}}}
  {\end{array}$\end{lrbox}%
   \ifdim\wd\progbox>\linewidth
     \PackageWarning{paper}{Overfull program display by \the\dimexpr\wd\progbox-\linewidth\relax}%
   \fi
   \[\usebox{\progbox}\]\ignorespacesafterend}

% written by cost.py; do not edit
\newcommand{\exBareHundred}{6.20}
\newcommand{\exBareTwoHundred}{22.40}
\newcommand{\exCases}{4,068}
\newcommand{\exCompShare}{46}
\newcommand{\exCostHundred}{0.78}
\newcommand{\exCostTwoHundred}{2.35}
\newcommand{\exCross}{62}
\newcommand{\exDouble}{3.0}
\newcommand{\exE}{10,900}
\newcommand{\exExact}{1,000}
\newcommand{\exFeasible}{190}
\newcommand{\exFirstExtra}{26,220}
\newcommand{\exFirstShare}{29}
\newcommand{\exPbest}{17}
\newcommand{\exPerTurnBest}{5,281}
\newcommand{\exPerTurnFull}{11,419}
\newcommand{\exPfull}{185}
\newcommand{\exPrograms}{3,000}
\newcommand{\exPstar}{16.5}
\newcommand{\exPstarBare}{6.3}
\newcommand{\exQuestion}{92,000}
\newcommand{\exRentShare}{74}
\newcommand{\exRhoFive}{0.26}
\newcommand{\exRhoOne}{1.00}
\newcommand{\exRhoTwenty}{0.13}
\newcommand{\exSlopeOne}{1.333}
\newcommand{\exSlopeTwo}{1.141}
\newcommand{\exSlopeZero}{2.000}
\newcommand{\exStepBest}{21}
\newcommand{\exStepE}{17,940}
\newcommand{\exStepPerTurn}{7,014}
\newcommand{\exStepStar}{21.2}
\newcommand{\exTabBareHundred}{6.20}
\newcommand{\exTabBareTenK}{50,120.00}
\newcommand{\exTabBareThousand}{512.00}
\newcommand{\exTabCompHundred}{0.54}
\newcommand{\exTabCompTenK}{52.82}
\newcommand{\exTabCompThousand}{5.31}
\newcommand{\exTabDelHundred}{0.62}
\newcommand{\exTabDelTenK}{68.85}
\newcommand{\exTabDelThousand}{6.17}
\newcommand{\exTabKHundred}{25}
\newcommand{\exTabKTenK}{47}
\newcommand{\exTabKThousand}{25}
\newcommand{\exTabPHundred}{20}
\newcommand{\exTabPTenK}{16}
\newcommand{\exTabPThousand}{17}
\newcommand{\exTabPlainHundred}{0.78}
\newcommand{\exTabPlainTenK}{4,037.21}
\newcommand{\exTabPlainThousand}{43.73}
\newcommand{\exTabStepHundred}{0.71}
\newcommand{\exTabStepTenK}{70.18}
\newcommand{\exTabStepThousand}{7.03}
\newcommand{\exTabWinHundred}{2.34}
\newcommand{\exTabWinTenK}{233.01}
\newcommand{\exTabWinThousand}{23.31}
\newcommand{\exWinContext}{255,000}
\newcommand{\exWinCost}{23,300}
\newcommand{\exWinLimit}{243,750}
\newcommand{\exWinTurn}{245}


\begin{document}

\title{What an Agent Loop Costs: Sized Types for Tokens}

\author{Claude-Fable-5.2}
\authornote{An AI model. It wrote this paper at the request of a human who asked for papers on
the strongest ideas of this repository, of at least twelve pages each, written to withstand
the scrutiny of a POPL review. The proofs are on paper and have not been checked by a machine.
A script tests the theorems on random programs. At the time of writing nobody has reviewed the
paper.}

\begin{abstract}
An agent that works in a loop sends its whole history to the model at every turn. The tokens it
pays for therefore grow with the square of the number of turns, and its author usually learns
this from the invoice. We give a type system that states the cost of a loop in advance. A
string type carries two sizes: a bound on a prefix that the provider has already stored, which
is billed at a reduced price, and a bound on the rest. A call to the model costs its prompt and
its answer, a loop is typed with an invariant on the sizes of its state, and an effect records
where the program may wait for so long that the provider's cache expires. We prove that the
cost in the type bounds the cost of every execution, whatever the model, the user, the tools
and the provider do, and that without a cache and without choice the bound is attained. Typing
six programs gives six laws. In a loop that only appends, every token is bought once and then
rented at every later turn. A question to a person forfeits the cache. A loop that summarizes
its history every $p$ turns has linear cost, which is least at $p = \sqrt{2E/(\rho d)}$. A
sliding window is cheaper than the full history only when it is smaller than the discounted
history. One token placed before the history costs as much as having no cache. With $L$ levels
of delegation, the overhead of $T$ turns grows as $T^{1+1/(2^{L+1}-1)}$. A script tests the
theorems on random programs and every formula against an implementation of the semantics.
\end{abstract}

\keywords{sized types, cost analysis, type-and-effect systems, language models, agents, caching}

\authorsaddresses{}

\maketitle

\section{Introduction}

An agent is a loop around a language model. At every turn the loop receives a message, from its
user or from a tool that the model asked to run, appends the message to a context, sends the
context to the model, and appends the answer~\cite{yao2022}. The model keeps no state between
calls, so every call carries the whole context, and the provider bills every call by the token:
a price for each token of the prompt and a higher price for each token of the answer. If a turn
adds $d$ tokens, then turn $k$ sends about $kd$ tokens, and $T$ turns send about $dT^2/2$. Twice
the turns cost up to four times as much.

The people who build agents know this in outline and answer it with engineering. Providers keep
a cache: a prompt that begins with a prefix sent a few minutes earlier is billed for that
prefix at a fraction $\rho$ of the price~\cite{anthropic2026}. Harnesses replace a long context
by a summary or keep only a window of recent messages~\cite{packer2023}, or they hand parts
of the work to agents with contexts of their own. Each of these measures changes the bill, and by how
much depends on details of the program: whether it appends to the context or writes at its
top, how often it summarizes, and where it waits for a person. Today these effects are found by
measurement~\cite{lumer2026,xu2026}, after the money is spent.

This paper treats the bill as a static property of the loop. A harness is a program, the price
list is public, and the quantities that decide the bill are the sizes of strings. We define a
small calculus of agent loops with a cost semantics, and a type system whose judgment
$\Gamma \vdash_\varphi e : A \mid C$ says that $e$ costs at most $C$. Three ideas fit the type
system to the price list.

\begin{itemize}
\item \emph{A string type has two sizes.} A value of type $\Str{P}{R}$ is a string that begins
with at most $P$ tokens which the provider has stored and continues with at most $R$ other
tokens. Sending it as a prompt costs at most $\rho P + R$. The split is the invariant of every
loop that appends: the context is what was sent last time, followed by the last answer and the
new message.
\item \emph{A call returns what is now stored.} The call $\call{v_1}{v_2}{M}$ sends the prompt
$v_1v_2$ and asks for an answer of at most $M$ tokens. It returns $v_1$, the prompt, and the
answer. The two returned strings are equal to strings that the program already had, but their
types are new: both are now known to be stored. The fact that a string was sent thus travels
with a value and needs no flow-sensitive environment.
\item \emph{Waiting is an effect.} A provider's cache forgets after a few minutes. The
expression $\hear{M}$ receives a message that may take long to arrive, such as an answer from a
person. Its effect $\varphi = 1$ makes the typing rules forget, for every variable in scope,
that a prefix was stored.
\end{itemize}

\paragraph{Results.} The cost semantics leaves four parties free: the model and the user choose
the content and the length of answers and messages within the stated bounds, a wait may empty
any part of the cache, and the provider may bill more of a prompt at the reduced price than it
must, never less. Theorem~\ref{thm:sound} states that the cost in the type bounds the cost of
every evaluation. Corollary~\ref{cor:window} states that no call of a typed program exceeds the
model's context window. Proposition~\ref{prop:attained} states that without a cache and without
choice the bound is attained, so that the rules lose nothing of their own.
Section~\ref{sec:laws} types six programs. Each typing derivation yields a closed form, and
each closed form is a law about agent loops; Table~\ref{tab:laws} lists them. Some of them are
folklore among practitioners, such as the advice to keep the stable part of a prompt in front.
The type system derives them from one set of rules, with constants, and makes them properties
of the program text that a checker can enforce.

\begin{table}[t]
\caption{The laws of Section~\ref{sec:laws}. A turn adds a message of at most $q$ tokens and an
answer of at most $m$ tokens, and $d = q + m$. The price of a stored prompt token is $\rho$,
that of an answer token is $\omega$, and the unit is the price of a prompt token that is not
stored. The system prompt has $h$ tokens.}
\label{tab:laws}
\small
\begin{tabular}{@{}lll@{}}
\toprule
Program & Cost of $T$ turns & Law \\
\midrule
plain loop, no cache (\S\ref{sec:plain}) & $dT^2/2 + O(T)$ & every turn buys the context again \\
plain loop (\S\ref{sec:plain}) & $\rho\,dT^2/2 + O(T)$ & a token is bought once, then rented \\
a person every $j$ turns (\S\ref{sec:human}) & $(\rho + \frac{1-\rho}{j})\,dT^2/2 + O(T)$ & a wait forfeits the cache \\
summary every $p$ turns (\S\ref{sec:compaction}) & $T\,(E/p + \rho dp/2 + c)$ & least at $p = \sqrt{2E/(\rho d)}$ \\
window of $b$ tokens (\S\ref{sec:window}) & $T\,(b + \rho h + q + \omega m)$ & pays only if $b < \rho \cdot \text{history} + m$ \\
window sliding every $p$ turns (\S\ref{sec:window}) & $T\,(E'/p + \rho dp/2 + c')$ & $E' = (1-\rho)(b - m)$ \\
a token before the history (\S\ref{sec:order}) & $dT^2/2 + O(T)$ & as without a cache \\
$L$ levels of delegation (\S\ref{sec:delegation}) & overhead $\Theta(T^{1+1/(2^{L+1}-1)})$ & $\theta_L - 1$ more than halves per level \\
\bottomrule
\end{tabular}
\end{table}

\paragraph{An example in numbers.} Throughout we use one example: a system prompt of
$h = 10{,}000$ tokens, messages and answers of at most $q = m = 500$ tokens, and the prices
$\rho = 0.08$ and $\omega = 4$, which Section~\ref{sec:calculus} takes from the price list of a
provider. A plain loop of 200 turns then costs \exBareTwoHundred{} million units without a
cache and \exCostTwoHundred{} million with one, and \exRentShare{} percent of the latter is
rent on tokens bought at earlier turns. Summarizing every \exPbest{} turns, the period that the
square-root law gives, costs \exCompShare{} percent of summarizing only when the context window
is full. A sliding window of 20,000 tokens costs three times as much as the plain loop over the
first 100 turns. All numbers in this paper are computed by the script of
Section~\ref{sec:tests}.

\paragraph{What the paper does not do.} We bound money and say nothing about the quality of
the agent's work: a summary forgets, and no term of our formulas prices that. We have not
compared the bounds with invoices; the semantics follows the published billing rules of one
provider, and Section~\ref{sec:calculus} states under which conditions. The calculus is first
order with bounded loops, which is the shape of a harness with a turn limit, and the
invariants of loops are given by the programmer. The proofs are on paper.
Section~\ref{sec:limits} lists these limits again, together with the problems that we think
come next.

\paragraph{Contributions.}
\begin{itemize}
\item A calculus of agent loops and a cost semantics with a prefix cache, parametric in the
prices and nondeterministic in the model, the user, and the provider
(Section~\ref{sec:calculus}).
\item A type-and-effect system with split string sizes, cost bounds and loop invariants
(Section~\ref{sec:types}), with a proof that its bounds hold, that they keep every call within
the context window, and that they are attained when there is no cache and no choice
(Section~\ref{sec:sound}).
\item Six programs with their types and the laws of Table~\ref{tab:laws}
(Section~\ref{sec:laws}), among them a square-root law for the period of compaction and the
exponent $1 + 1/(2^{L+1}-1)$ for delegation with $L$ levels.
\item An executable semantics and a checker for closed programs, used to test the theorems on
random programs and every closed form on a grid of parameters (Section~\ref{sec:tests}).
\end{itemize}

\section{Overview: The Plain Loop}
\label{sec:overview}

This section types the simplest agent and shows the two ways in which a program loses the
cache. The definitions are made precise in Sections~\ref{sec:calculus} and~\ref{sec:types}.

\paragraph{The loop.} One turn of an agent with context $y$ is the expression
\begin{prog}
\turn(y) \defeq \letin{u}{\recv{q}} & \text{a message of at most $q$ tokens}\\
\phantom{\turn(y) \defeq{}} \letin{(\_, c, a)}{\call{y}{u}{m}} & \text{send $y$ then $u$; the answer has at most $m$ tokens}\\
\phantom{\turn(y) \defeq{}} c \cat a & \text{the new context}
\end{prog}
The harness cuts every message and every tool result to $q$ tokens, and it passes $m$ to the
provider as the limit on the length of the answer. The call returns three strings: its first
argument, the prompt $c = yu$, and the answer $a$. An agent that starts from a system prompt
$w_0$ of $h$ tokens and runs for $N+1$ turns is
\[
\nm{loop}_N \defeq \letin{x}{\turn(w_0)}\; \iter{N}{x}{i, y}{\turn(y)} .
\]

\paragraph{Its type.} Suppose that $y$ has the type $\Str{P}{R}$: its first $P$ tokens are
stored at the provider and $R$ tokens follow. The prompt $yu$ then has $P$ stored tokens and
$R + q$ others, so the call costs at most $\rho P + R + q$ for the prompt and $\omega m$ for the
answer. After the call the whole prompt is stored, so $c$ has the type $\Str{P+R+q}{0}$, and
the new context $c \cat a$ has the type $\Str{P+R+q}{m}$. One turn thus moves everything that
was sent into the stored part and leaves the answer as the rest.

The literal $w_0$ has the type $\Str{0}{h}$, since nothing has been sent yet. The first turn
costs $h + q + \omega m$ and yields a context of type $\Str{h+q}{m}$. From there on the type of
the context after $i$ further turns is $\Str{h + q + i d}{m}$ with $d = q + m$: this is the
invariant of the loop, and turn $i$ of the loop costs $\rho\,(h + q + i d) + d + \omega m$.
Summing over $i < N$ and adding the first turn gives the cost of $\nm{loop}_N$:
\[
\underbrace{h + q + N d}_{\text{bought}} \;+\;
\underbrace{(N+1)\,\omega m}_{\text{answers}} \;+\;
\rho \underbrace{\tsum_{i<N} (h + q + i d)}_{\text{rented}} .
\]
Every token of the final context except the last answer is billed once at the full price, in
the first prompt that contains it. At every later turn it is billed again at the price $\rho$.
The first term is linear in $N$, and the last is quadratic. A cache therefore multiplies the
quadratic term by $\rho$ and leaves it quadratic.

\paragraph{Losing the cache by order.} Suppose that the harness writes the current time at the
top of every prompt, so that the call is $\call{t}{y \cat u}{m}$ for a new string $t$ of one
token. The rule for calls gives the reduced price only to the stored part of the first
argument. Here the first argument is new, the stored part of $y$ does not count because it no
longer stands at the front, and the turn costs $1 + P + R + q + \omega m$. No prompt of this
loop begins with a stored string, the context keeps the type $\Str{0}{h + id}$, and the loop
costs as much as it would without a cache (Corollary~\ref{cor:clock}). One token has multiplied
the rent by $1/\rho$.

\paragraph{Losing the cache by waiting.} Suppose that the message comes from a person, and
write $\hear{q}$ in place of $\recv{q}$. The person may answer after the cache has expired. The
rule for $\kw{let}$ therefore types what follows a $\kw{hear}$ in an environment that has
forgotten all stored prefixes: there $y$ has the type $\Str{0}{P+R}$ and the turn costs
$P + R + q + \omega m$. A question to a person costs $(1-\rho)P$ more than a message from a
tool (Corollary~\ref{cor:human}).

\paragraph{Why types.} The three analyses above are compositional: $\turn$ is typed once, as a
function of $P$ and $R$, and is used in five of the six programs of Section~\ref{sec:laws}. The
bound holds for every behavior of the model, which no test run can establish. And the
discipline that keeps a loop cheap, namely that a prompt begins with a string known to be
stored, is a property of the program text, so a type checker can report its violation before
the loop runs.

\section{A Calculus of Agent Loops}
\label{sec:calculus}

\paragraph{Strings and sizes.} A \emph{string} $w$ is a finite sequence of tokens, $|w|$ is
its length, $\eps$ is the empty string, and $u \pre w$ says that $u$ is a prefix of $w$. Sizes
and costs are written with index terms. An \emph{index term} denotes a natural number and a
\emph{cost term} a nonnegative real:
\[
I, J, M, N, P, R ::= i \mid n \mid I + J \mid I \cdot J \mid \tsum_{i<I} J
\qquad\qquad
C, D ::= I \mid r \cdot C \mid C + D \mid \tsum_{i<I} C
\]
Here $i$ is an index variable, $n$ a natural number, and $r$ a nonnegative real. A valuation
$\eta$ maps index variables to natural numbers, and $\sem{I}\eta$ and $\sem{C}\eta$ are the
values of the terms. An inequality $I \le J$ or $C \le D$ \emph{holds under} $\eta$ if
the values of its two sides under $\eta$ compare so.

\paragraph{Syntax.} Figure~\ref{fig:syntax} gives the syntax. Values are variables, strings and
pairs. The expressions $v_1 \cat v_2$ and $\tail{M}{v}$ concatenate strings and keep the last
$M$ tokens of a string. The expressions $\recv{M}$ and $\hear{M}$ receive a message of at most
$M$ tokens; the second one may take long. A message costs no tokens, and what a tool costs
otherwise is outside the model. The call $\call{v_1}{v_2}{M}$ sends the prompt
$v_1v_2$ to the model and returns a triple. The choice $e_1 \orr e_2$ runs one of its branches;
it stands for a test on the content of a string, about which the type system knows nothing,
and it lets a loop stop early by doing nothing in its remaining turns. The loop
$\iter{N}{v}{i,x}{e}$ runs $e$ exactly $N$ times, with $x$ bound to the state and $i$ to the
number of the turn. Every harness that we know has a limit on the number of turns, and $N$ is
that limit.

We use three abbreviations. A triple $(v_1, v_2, v_3)$ is the pair $(v_1, (v_2, v_3))$, and
$\kw{let}$ may bind a triple or a pair of an expression, which stands for the evident sequence
of bindings. A concatenation in the place of a value, as in $\call{y}{u \cat v}{M}$, stands for
a $\kw{let}$ that names it. The variable $\_$ is one that is not used.

\begin{figure}[t]
\[
\begin{array}{@{}l@{\quad}r@{\;}l@{}}
\text{values} & v ::= & x \mid w \mid (v_1, v_2) \\
\text{expressions} & e ::= & v \mid v_1 \cat v_2 \mid \tail{M}{v} \mid \recv{M} \mid \hear{M}
  \mid \call{v_1}{v_2}{M} \\
 & \mid & \letin{x}{e_1}\; e_2 \mid \letin{(x, y)}{v}\; e \mid e_1 \orr e_2
  \mid \iter{N}{v}{i, x}{e} \\
\text{types} & A, B ::= & \Str{P}{R} \mid A \times B
\end{array}
\]
\Description{The grammar of values, expressions and types.}
\caption{Syntax. $M$, $N$, $P$ and $R$ are index terms.}
\label{fig:syntax}
\end{figure}

\paragraph{Prices.} The unit of cost is the price of a prompt token that the provider does not
find in its cache. The price of a prompt token that it finds is $\rho$, with
$0 \le \rho \le 1$, and the price of a token of the answer is $\omega \ge 0$. For example, the
documentation of one provider~\cite{anthropic2026} lists for several of its models a base price
for prompt tokens, $1.25$ times the base price for prompt tokens that are written to the
cache, $0.1$ times for tokens read from the cache, and $5$ times for answer tokens. In a loop
all new prompt tokens are written to the cache, so our unit is the price of a written token,
and these multipliers give $\rho = 0.1/1.25 = 0.08$ and $\omega = 5/1.25 = 4$. Setting
$\rho = 1$ describes a provider without a cache.

\paragraph{The cache.} A \emph{cache} $K$ is a finite set of strings, the \emph{breakpoints}:
the prefixes of earlier prompts that the provider has stored. A call with the prompt $w_1w_2$
stores two breakpoints, $w_1$ and $w_1w_2$. The provider must bill at the reduced price the
longest breakpoint that the prompt begins with. We write
\[
\hit_K(w) = \max\,\{\, |u| \;:\; u \in K \cup \{\eps\},\; u \pre w \,\}
\]
for its length. A cache stores breakpoints and not all their prefixes: a prompt that shares
only half of a breakpoint with an earlier prompt need not be found.

\paragraph{Evaluation.} Figure~\ref{fig:semantics} defines the judgment
$\ev{K}{e}{c}{K'}{v}$: in the cache $K$ the closed expression $e$ evaluates to the value $v$
at the cost $c$ and leaves the cache $K'$. In a closed expression every index term is closed,
and we identify it with its value. The judgment is nondeterministic in four places. In
\textsc{E-Recv} and \textsc{E-Hear} the message is any string within the bound, chosen by a
user or a tool. In \textsc{E-Call} the answer is any string within the bound, chosen by the
model. Also in \textsc{E-Call}, the provider bills $\ell$ tokens at the reduced price, where
$\ell$ is any number between the length that it must grant and the length of the prompt. In
\textsc{E-Hear} the cache shrinks to any subset.

\begin{figure}[t]
\begin{mathpar}
\inferrule*[right=E-Val]{ }{\ev{K}{v}{0}{K}{v}}
\and
\inferrule*[right=E-Cat]{ }{\ev{K}{w_1 \cat w_2}{0}{K}{w_1w_2}}
\and
\inferrule*[right=E-Tail]{w = w'w'' \\ |w''| = \min(n, |w|)}{\ev{K}{\tail{n}{w}}{0}{K}{w''}}
\and
\inferrule*[right=E-Recv]{|w| \le n}{\ev{K}{\recv{n}}{0}{K}{w}}
\and
\inferrule*[right=E-Hear]{|w| \le n \\ K' \subseteq K}{\ev{K}{\hear{n}}{0}{K'}{w}}
\and
\inferrule*[right=E-Call]{|a| \le n \\ \hit_K(w_1w_2) \le \ell \le |w_1w_2|}
  {\ev{K}{\call{w_1}{w_2}{n}}{\rho\ell + (|w_1w_2| - \ell) + \omega|a|}{K \cup \{w_1, w_1w_2\}}{(w_1, w_1w_2, a)}}
\and
\inferrule*[right=E-Let]{\ev{K}{e_1}{c_1}{K_1}{v_1} \\ \ev{K_1}{e_2\{v_1/x\}}{c_2}{K_2}{v_2}}
  {\ev{K}{\letin{x}{e_1}\; e_2}{c_1 + c_2}{K_2}{v_2}}
\and
\inferrule*[right=E-Split]{\ev{K}{e\{v_1/x, v_2/y\}}{c}{K'}{v}}
  {\ev{K}{\letin{(x, y)}{(v_1, v_2)}\; e}{c}{K'}{v}}
\and
\inferrule*[right=E-Or]{\ev{K}{e_k}{c}{K'}{v} \\ k \in \{1, 2\}}{\ev{K}{e_1 \orr e_2}{c}{K'}{v}}
\and
\inferrule*[right=E-Iter0]{ }{\ev{K}{\iter{0}{v}{i,x}{e}}{0}{K}{v}}
\and
\inferrule*[right=E-Iter]{\ev{K}{\iter{n}{v}{i,x}{e}}{c_1}{K_1}{v_1} \\
  \ev{K_1}{e\{n/i, v_1/x\}}{c_2}{K_2}{v_2}}
  {\ev{K}{\iter{n+1}{v}{i,x}{e}}{c_1 + c_2}{K_2}{v_2}}
\end{mathpar}
\Description{Eleven inference rules that define the evaluation judgment.}
\caption{Cost semantics. $K$ is the cache, $c$ the cost, and $n$ a natural number.}
\label{fig:semantics}
\end{figure}

\paragraph{What the semantics assumes.} Three assumptions connect Figure~\ref{fig:semantics}
to a real provider. The first is that strings are sequences of tokens and that concatenation
adds lengths. A tokenizer may merge characters across the seam of two strings, and a chat
format adds a few tokens to every message; both change lengths by a constant for each
message, which can be added to $q$ and $m$. The second is that a prompt which begins with a
breakpoint is billed at the reduced price for at least that breakpoint. The documented cache
grants this when the breakpoint is longer than a minimum length and lies within a fixed number
of content blocks before a breakpoint of the new prompt~\cite{anthropic2026}; a turn that
adds more blocks than that needs an additional breakpoint. The third is that a breakpoint stays in the
cache until the next $\kw{hear}$. The documented cache keeps an entry for five minutes and
renews it at every use, so the programmer must write $\kw{hear}$ for every input that may take
longer and may write $\kw{recv}$ for the others. Evaluation always terminates, because a
loop runs for a fixed number of turns.

\section{The Type System}
\label{sec:types}

\paragraph{Types.} A type is a string type or a product (Figure~\ref{fig:syntax}). A value of
type $\Str{P}{R}$ is a string $w_1w_2$ such that $w_1$ is a breakpoint of the current cache or
empty, $|w_1| \le P$ and $|w_2| \le R$. We call $P$ the \emph{stored} size and $R$ the
\emph{rest}, and $P + R$ the total size. Both sizes are upper bounds. This is enough because
the cost of a prompt is monotone in both: a prompt with stored part $w_1$ and rest $w_2$ costs
at most $\rho|w_1| + |w_2|$, whatever else the provider finds.

\paragraph{Subtyping.} Two rules generate the subtype relation on string types, and products
are covariant:
\begin{mathpar}
\inferrule*[right=S-Widen]{P \le P' \\ R \le R'}{\Str{P}{R} \le \Str{P'}{R'}}
\and
\inferrule*[right=S-Forget]{ }{\Str{P}{R} \le \Str{0}{P + R}}
\and
\inferrule*[right=S-Pair]{A \le A' \\ B \le B'}{A \times B \le A' \times B'}
\end{mathpar}
The relation is closed under reflexivity and transitivity. A type may loosen its bounds, and it
may forget that a prefix is stored. It may not move a part of the stored size to the rest,
because a prefix of a breakpoint is not a breakpoint. The function $\fl{\cdot}$ forgets
everywhere: $\fl{\Str{P}{R}} = \Str{0}{P+R}$ and $\fl{A \times B} = \fl{A} \times \fl{B}$, and
$\fl{\Gamma}$ applies it to every type of an environment. We write $\Gamma^0$ for $\Gamma$ and
$\Gamma^1$ for $\fl{\Gamma}$.

\paragraph{Judgments.} Figure~\ref{fig:typing} defines two judgments. $\Gamma \vdash v : A$
types a value. $\Gamma \vdash_\varphi e : A \mid C$ says that $e$ evaluates to a value of type
$A$ at a cost of at most $C$; the effect $\varphi$ is $1$ if $e$ may wait and $0$ if it
cannot. Effects are ordered by $0 \le 1$, and $\vee$ is their maximum. A judgment is made under a set of \emph{assumptions}, which are
inequalities between index terms and which we do not write. The rule \textsc{Iter} adds the
assumption $i + 1 \le N$ for its body, and a program may come with assumptions on its
parameters. A premise $I \le J$ of a rule, also of a rule of subtyping, means that the
inequality holds under every valuation that satisfies the assumptions.

\begin{figure}[t]
\begin{mathpar}
\inferrule*[right=Var]{\Gamma(x) = A}{\Gamma \vdash x : A}
\and
\inferrule*[right=Lit]{ }{\Gamma \vdash w : \Str{0}{|w|}}
\and
\inferrule*[right=Pair]{\Gamma \vdash v_1 : A_1 \\ \Gamma \vdash v_2 : A_2}{\Gamma \vdash (v_1, v_2) : A_1 \times A_2}
\and
\inferrule*[right=V-Sub]{\Gamma \vdash v : A \\ A \le A'}{\Gamma \vdash v : A'}
\and
\inferrule*[right=Val]{\Gamma \vdash v : A}{\Gamma \vdash_0 v : A \mid 0}
\and
\inferrule*[right=Recv]{ }{\Gamma \vdash_0 \recv{M} : \Str{0}{M} \mid 0}
\and
\inferrule*[right=Hear]{ }{\Gamma \vdash_1 \hear{M} : \Str{0}{M} \mid 0}
\and
\inferrule*[right=Cat]{\Gamma \vdash v_1 : \Str{P_1}{R_1} \\ \Gamma \vdash v_2 : \Str{P_2}{R_2}}
  {\Gamma \vdash_0 v_1 \cat v_2 : \Str{P_1}{R_1 + P_2 + R_2} \mid 0}
\and
\inferrule*[right=Tail]{\Gamma \vdash v : \Str{P}{R}}{\Gamma \vdash_0 \tail{M}{v} : \Str{0}{M} \mid 0}
\and
\inferrule*[right=Call]{\Gamma \vdash v_1 : \Str{P_1}{R_1} \\ \Gamma \vdash v_2 : \Str{P_2}{R_2} \\ I = P_1 + R_1 + P_2 + R_2}
  {\Gamma \vdash_0 \call{v_1}{v_2}{M} : \Str{P_1 + R_1}{0} \times \Str{I}{0} \times \Str{0}{M}
   \mid \rho P_1 + R_1 + P_2 + R_2 + \omega M}
\and
\inferrule*[right=Let]{\Gamma \vdash_{\varphi_1} e_1 : A \mid C_1 \\ \Gamma^{\varphi_1}, x : A \vdash_{\varphi_2} e_2 : B \mid C_2}
  {\Gamma \vdash_{\varphi_1 \vee \varphi_2} \letin{x}{e_1}\; e_2 : B \mid C_1 + C_2}
\and
\inferrule*[right=Split]{\Gamma \vdash v : A_1 \times A_2 \\ \Gamma, x : A_1, y : A_2 \vdash_\varphi e : B \mid C}
  {\Gamma \vdash_\varphi \letin{(x, y)}{v}\; e : B \mid C}
\and
\inferrule*[right=Or]{\Gamma \vdash_\varphi e_1 : A \mid C \\ \Gamma \vdash_\varphi e_2 : A \mid C}
  {\Gamma \vdash_\varphi e_1 \orr e_2 : A \mid C}
\and
\inferrule*[right=Iter]{\Gamma \vdash v : A\{0/i\} \\ \Gamma^{\varphi}, x : A \vdash_\varphi e : A\{i+1/i\} \mid C \\ i \notin \mathrm{fv}(\Gamma, N)}
  {\Gamma \vdash_\varphi \iter{N}{v}{i, x}{e} : A\{N/i\} \mid \tsum_{i<N} C}
\and
\inferrule*[right=Sub]{\Gamma \vdash_\varphi e : A \mid C \\ A \le A' \\ C \le C' \\ \varphi \le \varphi'}
  {\Gamma \vdash_{\varphi'} e : A' \mid C'}
\end{mathpar}
\Description{Sixteen inference rules that define the two typing judgments.}
\caption{Typing rules for values and expressions.}
\label{fig:typing}
\end{figure}

\paragraph{Strings.} A literal has nothing stored (\textsc{Lit}), and neither has a message
(\textsc{Recv}, \textsc{Hear}). A concatenation keeps the stored part of its left argument
only (\textsc{Cat}): the stored part of the right argument no longer stands at the front. The
last $M$ tokens of a string are not a prefix of anything that was sent, so nothing of them is
stored (\textsc{Tail}).

\paragraph{Calls.} The rule \textsc{Call} is the center of the system. The prompt is $v_1v_2$.
Of its tokens only the stored part of $v_1$ is billed at $\rho$; the rest of $v_1$ and all of
$v_2$ are billed in full, and the answer costs $\omega M$. The call stores two breakpoints, and
its result type says so: the first component is $v_1$ and the second is the prompt, both with
their total size as stored size. The two-part form lets a program keep a stable front, such as
a system prompt, as a stored string of its own, which it can put in front of a new context
later. Section~\ref{sec:compaction} uses this after a summary. A call with a single prompt
$v$ is $\call{v}{\eps}{M}$.

An alternative design would update the type of a variable in the environment after it has been
sent, as a typestate. Returning the stored strings as new values avoids this: environments
stay flow-insensitive, and a variable that was bound before the call keeps its old, weaker
type, which remains valid.

\paragraph{Waiting.} The effect of $\kw{hear}$ is $1$. In \textsc{Let}, when the first
expression may wait, the second is typed in $\fl{\Gamma}$: the variables bound before the wait
keep their sizes and lose their stored parts. The variable $x$ that the $\kw{let}$ binds keeps
its type $A$, which was established after the wait. The rule \textsc{Iter} does the same for a
loop whose body may wait, since the body runs after earlier turns.

\paragraph{Loops.} In \textsc{Iter} the type $A$ may mention the index variable $i$, and it is
the invariant of the loop: the state has the type $A\{0/i\}$ at the start, the body takes
$A$ to $A\{i+1/i\}$, and the result has the type $A\{N/i\}$. The cost $C$ of turn $i$ may
mention $i$ as well, and the loop costs their sum.

A loop that starts from a new context needs some care. Its first turn differs from the
others in its type: before it nothing is stored, and after every turn almost everything is.
An invariant that covered both would need a case distinction on $i$, which index terms lack.
We therefore write the first turn before the loop, as in $\nm{loop}_N$ of
Section~\ref{sec:overview}.

\section{The Bounds Hold}
\label{sec:sound}

\paragraph{Interpretation.} For closed types and a cache $K$ we define the set of closed
values of a type:
\begin{align*}
\sem{\Str{P}{R}}_K &= \{\, w_1w_2 \mid w_1 \in K \cup \{\eps\},\; |w_1| \le P,\; |w_2| \le R \,\} \\
\sem{A \times B}_K &= \{\, (v_1, v_2) \mid v_1 \in \sem{A}_K,\; v_2 \in \sem{B}_K \,\}
\end{align*}
A substitution $\gamma$ of closed values for variables is in $\sem{\Gamma}_K$ if
$\gamma(x) \in \sem{\Gamma(x)}_K$ for every $x$. We write $A\eta$ for $A$ with its index
variables replaced by their values under $\eta$.

\begin{lemma}[Growth]
\label{lem:growth}
If $K \subseteq K'$ then $\sem{A}_K \subseteq \sem{A}_{K'}$.
\end{lemma}

\begin{lemma}[Forgetting]
\label{lem:forget}
$\sem{A}_K \subseteq \sem{\fl{A}}_{K'}$ for all caches $K$ and $K'$.
\end{lemma}

\begin{lemma}[Subtyping]
\label{lem:sub}
If $A \le A'$ then $\sem{A\eta}_K \subseteq \sem{A'\eta}_K$ for every cache $K$ and every
$\eta$ that satisfies the assumptions.
\end{lemma}

\begin{lemma}[Values]
\label{lem:values}
If $\Gamma \vdash v : A$ and $\gamma \in \sem{\Gamma\eta}_K$ then $v\gamma \in \sem{A\eta}_K$.
\end{lemma}

\begin{proof}
Lemma~\ref{lem:growth} holds because the cache occurs in the definition only in the condition
$w_1 \in K \cup \{\eps\}$. For Lemma~\ref{lem:forget}, a string $w_1w_2$ with $|w_1| \le P$ and
$|w_2| \le R$ is $\eps$ followed by a string of at most $P + R$ tokens, and $\eps$ needs no
cache. Lemma~\ref{lem:sub} is by induction on the derivation of $A \le A'$: \textsc{S-Widen}
keeps the same split, and \textsc{S-Forget} is the previous argument. Lemma~\ref{lem:values}
is by induction on the derivation, with Lemma~\ref{lem:sub} for \textsc{V-Sub}; a literal $w$
is $\eps$ followed by $w$.
\end{proof}

\begin{theorem}[The bound holds]
\label{thm:sound}
Let $\Gamma \vdash_\varphi e : A \mid C$. Let $\eta$ be a valuation that
satisfies the assumptions of the judgment, $K$ a cache, and
$\gamma \in \sem{\Gamma\eta}_K$. If $\ev{K}{e\eta\gamma}{c}{K'}{v}$, then
\begin{enumerate}
\item $v \in \sem{A\eta}_{K'}$,
\item $c \le \sem{C}\eta$, and
\item if $\varphi = 0$ then $K \subseteq K'$.
\end{enumerate}
\end{theorem}

\begin{proof}
By induction on the typing derivation. We omit $\eta$ where it only evaluates index terms.

\textsc{Val}: Lemma~\ref{lem:values}. \textsc{Recv}, \textsc{Tail}: the result has at most $M$
tokens and is $\eps$ followed by itself; the cache is unchanged. \textsc{Hear}: the same for
the result, and nothing is claimed about $K'$ since $\varphi = 1$.

\textsc{Cat}: By Lemma~\ref{lem:values} the arguments are strings $w_1 = uu'$ with
$u \in K \cup \{\eps\}$, $|u| \le P_1$, $|u'| \le R_1$, and $w_2$ with
$|w_2| \le P_2 + R_2$. The result $w_1w_2$ is $u$ followed by $u'w_2$, which has at most
$R_1 + P_2 + R_2$ tokens.

\textsc{Call}: With $w_1 = uu'$ and $w_2$ as in the previous case, let $w = w_1w_2$. Since
$u \in K \cup \{\eps\}$ and $u \pre w$, we have $\hit_K(w) \ge |u|$, so the number $\ell$ of
\textsc{E-Call} satisfies $\ell \ge |u|$. Because $\rho \le 1$,
\begin{align*}
c = |w| - (1-\rho)\,\ell + \omega|a| &\le |w| - (1-\rho)\,|u| + \omega|a| \\
  &= \rho|u| + |u'| + |w_2| + \omega|a| \;\le\; \rho P_1 + R_1 + P_2 + R_2 + \omega M .
\end{align*}
The new cache $K' = K \cup \{w_1, w\}$ contains $K$. It contains $w_1$, and
$|w_1| \le P_1 + R_1$, so $w_1 \in \sem{\Str{P_1+R_1}{0}}_{K'}$; likewise for $w$. The answer
has at most $M$ tokens.

\textsc{Let}: The evaluation ends with \textsc{E-Let}. The induction hypothesis for $e_1$
gives $v_1 \in \sem{A}_{K_1}$ and $c_1 \le C_1$. If $\varphi_1 = 0$ it also gives
$K \subseteq K_1$, so $\gamma \in \sem{\Gamma}_{K_1}$ by Lemma~\ref{lem:growth}. If
$\varphi_1 = 1$ then $\gamma \in \sem{\fl{\Gamma}}_{K_1}$ by Lemma~\ref{lem:forget}. In both
cases $\gamma[x \mapsto v_1] \in \sem{\Gamma^{\varphi_1}, x : A}_{K_1}$, and the induction
hypothesis for $e_2$ gives $v_2 \in \sem{B}_{K_2}$ and $c_2 \le C_2$. If
$\varphi_1 \vee \varphi_2 = 0$ then $K \subseteq K_1 \subseteq K_2$.

\textsc{Split}, \textsc{Or}: directly by the induction hypothesis.

\textsc{Iter}: Let $\eta_n = \eta[i \mapsto n]$. For $n < \sem{N}\eta$ the valuation
$\eta_n$ satisfies the assumption $i + 1 \le N$ of the body. We show by induction on
$n \le \sem{N}\eta$ that every evaluation $\ev{K}{(\iter{n}{v}{i,x}{e})\eta\gamma}{c}{K'}{v'}$ has
$v' \in \sem{A\eta_n}_{K'}$, $c \le \sum_{k<n} \sem{C}\eta_k$, and $K \subseteq K'$ if
$\varphi = 0$. For $n = 0$ this is Lemma~\ref{lem:values} for the first premise. For $n + 1$
the evaluation ends with \textsc{E-Iter}. The inner hypothesis gives
$v_1 \in \sem{A\eta_n}_{K_1}$, the bound on $c_1$, and $K \subseteq K_1$ if $\varphi = 0$. As
in the case of \textsc{Let}, $\gamma \in \sem{\Gamma^\varphi\eta}_{K_1}$, by
Lemma~\ref{lem:growth} if $\varphi = 0$ and by Lemma~\ref{lem:forget} otherwise. Since $i$ is
not free in $\Gamma$, the outer induction hypothesis applies to the body with the valuation
$\eta_n$, the cache $K_1$ and the substitution $\gamma[x \mapsto v_1]$. It gives
$v_2 \in \sem{A\{i+1/i\}\eta_n}_{K_2} = \sem{A\eta_{n+1}}_{K_2}$ and
$c_2 \le \sem{C}\eta_n$, and $K_1 \subseteq K_2$ if $\varphi = 0$. The claim of the theorem
is the case $n = \sem{N}\eta$, because $A\{N/i\}\eta = A\eta_n$ and
$\sem{\sum_{i<N} C}\eta = \sum_{k<n} \sem{C}\eta_k$.

\textsc{Sub}: Lemma~\ref{lem:sub} and the validity of $C \le C'$; the third claim is weaker
for $\varphi' = 1$.
\end{proof}

\begin{corollary}[Closed programs]
\label{cor:closed}
If $\vdash_\varphi e : A \mid C$ for a closed expression $e$ and closed $A$ and $C$, then
every evaluation of $e$ from the empty cache costs at most $C$.
\end{corollary}

The corollary quantifies over all evaluations. It therefore holds for every model, for every
sequence of messages, for every provider that grants at least the hits of
Section~\ref{sec:calculus}, and for every pattern of expiry at the points marked by
$\kw{hear}$.

\begin{corollary}[The context window]
\label{cor:window}
Let $W$ be a natural number. Suppose that in a derivation of $\vdash_\varphi e : A \mid C$
for closed $e$ every use of \textsc{Call} satisfies $P_1 + R_1 + P_2 + R_2 + M \le W$ under
its assumptions. Then in every
evaluation of $e$ from the empty cache, every call has a prompt and a limit on its answer that
together have at most $W$ tokens.
\end{corollary}

\begin{proof}
By the same induction: in the case of \textsc{Call} the prompt $w$ has at most
$P_1 + R_1 + P_2 + R_2$ tokens.
\end{proof}

A provider rejects a call that does not fit the context window of its model. A harness
usually finds this out at the turn where it happens. With the side condition of the corollary
the type checker finds it before the first turn.

\paragraph{Tightness.} The rules can lose precision in four places: at a choice, where one
branch may be cheaper; at a use of subsumption that loosens a bound; at a $\kw{tail}$ of a
string shorter than the window; and at the cache, when the provider finds more than the types
know. The next proposition says that there is no other place. The total size $|A|$ of a type
is the sum of $P + R$ over its string types. A derivation is \emph{exact} if it does not use
\textsc{Or}, every use of \textsc{Tail} has $M \le P + R$, and every use of \textsc{Sub} and
\textsc{V-Sub} relates types of equal total size, component by component, and equal costs,
where equal means equal under every valuation that satisfies the assumptions. An
evaluation is \emph{maximal} if every message and every answer in it has exactly as many
tokens as its bound allows.

\begin{proposition}[The bound is attained]
\label{prop:attained}
Let $\rho = 1$. If a closed expression $e$ has an exact derivation of
$\vdash_\varphi e : A \mid C$, then every maximal evaluation of $e$ from the empty cache costs
exactly $C$.
\end{proposition}

\begin{proof}
We prove a stronger statement for open expressions: if every string of $\gamma$ has exactly
the total size of its type in $\Gamma\eta$, then every maximal evaluation of $e\eta\gamma$
yields a value whose strings have exactly the total sizes of $A\eta$, at the cost
$\sem{C}\eta$. The proof is by induction on the exact derivation. A literal has its exact
size. \textsc{Cat} adds sizes. In \textsc{Tail} the string has $P + R \ge M$ tokens, so its
last $M$ tokens are $M$ tokens. In \textsc{Call} the cost is $|w| + \omega|a|$ for every
$\ell$ since $\rho = 1$, and this is $P_1 + R_1 + P_2 + R_2 + \omega M$. \textsc{S-Forget}
keeps total sizes. The other cases follow from the induction hypothesis.
\end{proof}

With a cache the bound can be loose for the fourth reason only when the program builds a
string a second time that was already sent, or when two prompts happen to begin alike; the
provider then finds a breakpoint that no type records. Section~\ref{sec:tests} reports for
each program of Section~\ref{sec:laws} whether its bound is attained for $\rho < 1$.

\begin{remark}[Polynomial cost]
Every index term and every cost term denotes a polynomial in its variables, since the sum of a
polynomial over $i < N$ is a polynomial in $N$. A program with a type therefore has a cost
that is polynomial in its loop bounds. The loop $\iter{N}{w}{i,x}{x \cat x}$, whose context
doubles at every turn, has no type unless $w$ is empty: the total size of an invariant would
have to double with $i$.
\end{remark}

\begin{remark}[Checking]
Checking a derivation with given invariants reduces to inequalities between index terms. Since
the terms denote polynomials, equalities are decided by normal forms. Inequalities between
polynomials over the natural numbers are undecidable in general~\cite{davis1973}; comparing
coefficients is a sound test. The side conditions of the programs in this paper are
equalities that hold without assumptions, apart from those of Corollary~\ref{cor:window}. For closed programs no invariant is needed at all: a checker can unroll the
loops and compute with numbers, and this is what the script of Section~\ref{sec:tests} does.
\end{remark}

\section{Six Laws}
\label{sec:laws}

We now type six programs. In all of them a turn receives a message of at most $q$ tokens and
an answer of at most $m$ tokens, and $d = q + m$ is the growth of the context per turn. The
letters $h$, $q$, $m$, $s$, $g$, $b$, $t$ and $r$ stand for natural numbers, and so do the
periods $j$, $p$, $k$ and $J$, which are at least $1$, so that $p - 1$ is a number. The loop
bounds $N$, $T$ and $n$ are index variables. The example uses $h = 10{,}000$, $q = m = 500$,
$\rho = 0.08$ and $\omega = 4$.

\begin{lemma}[A turn]
\label{lem:turn}
Let $\turn$ be the expression of Section~\ref{sec:overview}, and let $\hturn$ be the same
expression with $\hear{q}$ in place of $\recv{q}$. If $\Gamma \vdash v : \Str{P}{R}$ then
\begin{align*}
\Gamma &\vdash_0 \turn(v) : \Str{P + R + q}{m} \mid \rho P + R + q + \omega m, \\
\Gamma &\vdash_1 \hturn(v) : \Str{P + R + q}{m} \mid P + R + q + \omega m .
\end{align*}
\end{lemma}

\begin{proof}
For $\turn$: $u$ has the type $\Str{0}{q}$ by \textsc{Recv}, and \textsc{Call} gives the cost
$\rho P + R + 0 + q + \omega m$ and the types $c : \Str{P+R+q}{0}$ and $a : \Str{0}{m}$, so
$c \cat a$ has the type $\Str{P+R+q}{m}$ by \textsc{Cat}. For $\hturn$ the first binding has
the effect $1$, so \textsc{Let} types the rest in $\fl{\Gamma}$, where $v$ has the type
$\Str{0}{P+R}$: forgetting commutes with the rules for values. The call then costs
$P + R + q + \omega m$, and the types of $c$ and $a$ are as before.
\end{proof}

\subsection{The Plain Loop: Bought Once, Rented Every Turn}
\label{sec:plain}

\begin{corollary}
\label{cor:plain}
Let $\nm{loop}_N$ be the program of Section~\ref{sec:overview}, which runs one turn on a
string $w_0$ of $h$ tokens and then $N$ more turns. Then
\[
\vdash_0 \nm{loop}_N : \Str{h + q + Nd}{m} \;\mid\;
(h + q + Nd) + (N+1)\,\omega m + \rho\,\bigl(N(h+q) + d\,N(N-1)/2\bigr) .
\]
\end{corollary}

\begin{proof}
By Lemma~\ref{lem:turn} the first turn has the type $\Str{h+q}{m}$ and costs
$h + q + \omega m$. For the loop take the invariant $\Str{h + q + id}{m}$. By the lemma its
body has the type $\Str{h + q + id + m + q}{m}$, which is the invariant at $i + 1$, and costs
$\rho\,(h + q + id) + d + \omega m$. The sum over $i < N$ is
$\rho\,(N(h+q) + d\,N(N-1)/2) + Nd + N\omega m$.
\end{proof}

The three summands are the tokens bought, the answers, and the rent
(Section~\ref{sec:overview}). For a provider without a cache, $\rho = 1$ and the cost of
$T = N + 1$ turns is $T\,(h + q + \omega m) + d\,T(T-1)/2$: turn $k$ sends $h + kd + q$ tokens.
With a cache the quadratic term is multiplied by $\rho$. In the example, 100 turns cost
\exBareHundred{} million units without a cache and \exCostHundred{} million with one; 200 turns
cost \exBareTwoHundred{} and \exCostTwoHundred{} million. With the cache, twice the turns
still cost \exDouble{} times as much. In a loop of \exCross{} turns or more the rent is more
than half of the bill, and at 200 turns it is \exRentShare{} percent. The cache postpones the
quadratic regime and does not remove it. The largest prompt of the loop, together with the limit on its answer,
has $h + (N+1)d$ tokens. By Corollary~\ref{cor:window}, under the assumption
$h + (N+1)d \le W$ no call exceeds a context window of $W$ tokens; for $W = 200{,}000$ this
allows \exFeasible{} turns.

\subsection{A Person in the Loop: A Wait Forfeits the Cache}
\label{sec:human}

Suppose that every $j$-th message comes from a person and the others from tools:
\[
\nm{session}_n \defeq \iter{n}{w_0}{l, y}{\;\letin{x}{\hturn(y)}\; \iter{j-1}{x}{i,z}{\turn(z)}} .
\]

\begin{corollary}
\label{cor:human}
With $w_0$ of $h$ tokens, $\vdash_1 \nm{session}_n : \Str{0}{h + njd} \mid C$ where
\[
C = \bigl(1 + \rho(j-1)\bigr)\Bigl(nh + jd\,\frac{n(n-1)}{2}\Bigr)
  + n\Bigl(q + (j-1)d + j\omega m + \rho\bigl((j-1)q + d\,\frac{(j-1)(j-2)}{2}\bigr)\Bigr) .
\]
\end{corollary}

\begin{proof}
Take the invariant $\Str{0}{H}$ with $H = h + ljd$ for the outer loop. By
Lemma~\ref{lem:turn}, $\hturn(y)$ has the type $\Str{H+q}{m}$ and costs
$H + q + \omega m$. The inner loop has the invariant $\Str{H + q + id}{m}$ and costs
$\sum_{i<j-1} (\rho\,(H + q + id) + d + \omega m)$. Its result has the total size
$H + jd$, and \textsc{S-Forget} gives it the type $\Str{0}{H + jd}$, the invariant at
$l + 1$. Nothing is lost by this step, since the next turn waits. One block costs
$(1 + \rho(j-1))\,H + q + (j-1)d + j\omega m + \rho\,((j-1)q + d\,(j-1)(j-2)/2)$, and
$\sum_{l<n} H = nh + jd\,n(n-1)/2$.
\end{proof}

In terms of the number of turns $T = nj$, the quadratic term of $C$ is
\[
\Bigl(\rho + \frac{1-\rho}{j}\Bigr)\,\frac{d\,T^2}{2} .
\]
The cache thus acts with the weaker discount $\rho_j = \rho + (1-\rho)/j$. For $j = 1$, a
chat in which the person answers every turn, $\rho_1 = 1$ and the cache does nothing. In the
example $\rho_5 = \exRhoFive$ and $\rho_{20} = \exRhoTwenty$, against $\rho = 0.08$. Seen from
one turn, a question to a person costs $(1-\rho)P$ more than a message from a tool, where $P$
is the stored size of the context: with $P = 100{,}000$ this is \exQuestion{} units, the price
of 23,000 tokens of answer.

The type system is pessimistic here: $\kw{hear}$ stands for an input that may arrive after the
cache has expired, and a person who answers at once costs no more than a tool. But the worst
case is the common one for approvals. An agent that stops to ask for permission and is
answered after lunch pays for its whole context again. An earlier paper in this series argued
that an agent should ask only when an action is influenced by untrusted
data~\cite{tpp5}, because attention is scarce. The corollary adds a second reason, in tokens.

\subsection{Compaction: A Square-Root Law}
\label{sec:compaction}

A compacting agent replaces its context by a summary every $p$ turns. Its state is a pair of
the system prompt $\sigma$, which stays, and the current summary $z$. Let $\kappa$ be a
literal of $g$ tokens that asks for a summary of at most $s$ tokens.
\begin{prog}
\nm{block}(\sigma, z) \defeq \letin{u}{\recv{q}} & u : \Str{0}{q} \\
\quad \letin{(\sigma', c, a)}{\call{\sigma}{z \cat u}{m}} & \sigma' : \Str{h}{0},\; c : \Str{h+s+q}{0} \\
\quad \letin{y}{\iter{p-1}{c \cat a}{i, y}{\turn(y)}} & y : \Str{h + s + q + (p-1)d}{m} \\
\quad \letin{(\_, \_, z')}{\call{y}{\kappa}{s}} & z' : \Str{0}{s} \\
\quad (\sigma', z')
\end{prog}
The first call sends the system prompt as a part of its own, so that it comes back as a stored
string. The call for the summary puts the instruction $\kappa$ after the history.

\begin{corollary}
\label{cor:block}
If $\Gamma \vdash \sigma : \Str{h}{0}$ and $\Gamma \vdash z : \Str{0}{s}$, then
$\Gamma \vdash_0 \nm{block}(\sigma, z) : \Str{h}{0} \times \Str{0}{s} \mid U(p)$ with
\[
U(p) = (s + pd + g) + \omega\,(pm + s) + \rho\,\bigl(h + p\,(h + s + q) + d\,p(p-1)/2\bigr) .
\]
If $\sigma$ has the type $\Str{0}{h}$ instead, the cost is $U(p) + (1-\rho)h$. Hence the
program that runs one block on a literal system prompt and $n$ more blocks in a loop costs
$(n+1)\,U(p) + (1-\rho)h$.
\end{corollary}

\begin{proof}
The first call costs $\rho h + s + q + \omega m$. The loop has the invariant
$\Str{h + s + q + id}{m}$ and costs
$\sum_{i<p-1} (\rho\,(h + s + q + id) + d + \omega m)$ by Lemma~\ref{lem:turn}. The last call
costs $\rho\,(h + s + q + (p-1)d) + m + g + \omega s$. In the sum, the terms without $\rho$
and $\omega$ are $s + q + (p-1)d + m + g = s + pd + g$, the terms with $\omega$ are
$\omega\,(pm + s)$, and the terms with $\rho$ are
$\rho\,(h + \sum_{i<p} (h + s + q + id))$. If $\sigma$ is not stored, the first call costs
$h$ in place of $\rho h$.
\end{proof}

The cost per turn has the form
\[
\frac{U(p)}{p} = \frac{E}{p} + \frac{\rho d}{2}\,p + c_0,
\qquad E = (1+\omega)\,s + g + \rho h,
\qquad c_0 = d + \omega m + \rho\,(h + s + q) - \rho d/2 .
\]
It is linear in the number of turns for every $p$. The fixed cost $E$ of a compaction is
spread over $p$ turns, and the rent grows with $p$. $E$ consists of the summary, which is
bought once from the model at the price $\omega$ and once more as a part of the next prompt,
of the instruction, and of one more reading of the system prompt.

\begin{corollary}[Square-root law]
\label{cor:sqrt}
For $\rho > 0$ the cost per turn is least at $p^* = \sqrt{2E/(\rho d)}$, where it is
$\sqrt{2E\rho d} + c_0$. Over the natural numbers it is least at one of the two neighbors of
$p^*$.
\end{corollary}

\begin{proof}
The function $E/p + \rho dp/2$ is convex on $p > 0$ and its derivative vanishes at $p^*$.
\end{proof}

This is the formula for the best interval between checkpoints of a long
computation~\cite{young1974}, with the cost of a compaction in the place of the cost of a
checkpoint and the rent $\rho d$ in the place of the failure rate. In the example, with
$s = 2{,}000$ and $g = 100$, we get $E = \text{\exE}$ and $p^* = \exPstar$; the best whole period is
\exPbest{} turns, at \exPerTurnBest{} units per turn. A harness that compacts only when the
window of 200,000 tokens is full has $p = \exPfull$ and pays \exPerTurnFull{} units per turn.
Compacting at the best period costs \exCompShare{} percent of that. Without a cache the best
period is $\exPstarBare$ turns: the cache makes compaction less urgent.

Two cautions apply. The law minimizes money. A summary loses information, later turns may be
worse for it, and the formula has no term for that; it says what each compaction costs, so
that its price can be weighed against its damage. And the corollary assumes that turns have
equal bounds. For turns of different sizes the same derivation gives a sum in place of
$pd$, and the best moment to compact depends on the tokens accumulated, not on the number of
turns.

\subsection{Sliding Windows}
\label{sec:window}

A harness with a sliding window sends the system prompt and the last $b$ tokens of the
history. Its state is the pair of the stored system prompt and the window $x$:
\begin{prog}
\nm{wturn}(\sigma, x) \defeq \letin{u}{\recv{q}} \\
\quad \letin{(\sigma', \_, a)}{\call{\sigma}{x \cat u}{m}} & \text{cost } \rho h + b + q + \omega m \\
\quad \letin{x'}{\tail{b}{(x \cat u \cat a)}}\; (\sigma', x') & x' : \Str{0}{b}
\end{prog}

\begin{corollary}
\label{cor:window-cost}
If $\Gamma \vdash \sigma : \Str{h}{0}$ and $\Gamma \vdash x : \Str{0}{b}$ then
$\Gamma \vdash_0 \nm{wturn}(\sigma, x) : \Str{h}{0} \times \Str{0}{b} \mid \rho h + b + q + \omega m$.
A first turn on a literal system prompt followed by $N$ turns in a loop costs
$(h + b + q + \omega m) + N\,(\rho h + b + q + \omega m)$.
\end{corollary}

\begin{proof}
\textsc{Call} with $\sigma : \Str{h}{0}$ and $x \cat u : \Str{0}{b+q}$, then \textsc{Tail}.
The invariant of the loop is the type of the state, which does not depend on $i$.
\end{proof}

The cost is linear in the number of turns, and every turn pays the full price for the window:
the window of one turn is not a prefix of anything sent before, because its first tokens were
dropped. Compare a turn of the plain loop whose context has the stored size $P$, which costs
$\rho P + d + \omega m$ by Lemma~\ref{lem:turn}. The plain turn is the cheaper one exactly when
\[
\rho\,(P - h) + m < b ,
\]
that is, when the stored history without the system prompt is shorter than $(b - m)/\rho$. A
window saves money only if it is smaller than a fraction $\rho$ of the history that it
replaces, plus one answer. In the example a window of $b = 20{,}000$ tokens costs \exWinCost{} units per turn,
and the plain loop reaches that price per turn only at turn \exWinTurn, with a context of
\exWinContext{} tokens, which no longer fits the context window. Up to the limit of the
model, keeping everything is cheaper than this window at every single turn.

The loss comes from sliding at every turn. A window that slides in steps keeps the cache
between the steps. The following block runs $p$ turns on a growing context and then keeps the
last $b$ tokens of the history, which it carries beside the context as $k$:
\begin{prog}
\nm{sblock}(\sigma, x) \defeq \letin{u}{\recv{q}} \\
\quad \letin{(\sigma', c, a)}{\call{\sigma}{x \cat u}{m}} & c : \Str{h + b + q}{0} \\
\quad \letx{(y, k)}{\kw{iter}\;p-1\;\kw{from}\;(c \cat a,\; x \cat u \cat a)\;\kw{as}\;i, (y, k).} \\
\quad\qquad \letin{u}{\recv{q}}\; \letin{(\_, c, a)}{\call{y}{u}{m}} & y : \Str{h + b + q + id}{m} \\
\quad\qquad (c \cat a,\; k \cat u \cat a)\;\kw{in} & k : \Str{0}{b + d + id} \\
\quad \letin{x'}{\tail{b}{k}}\; (\sigma', x')
\end{prog}

\begin{corollary}
\label{cor:step}
If $\Gamma \vdash \sigma : \Str{h}{0}$ and $\Gamma \vdash x : \Str{0}{b}$, then
$\Gamma \vdash_0 \nm{sblock}(\sigma, x) : \Str{h}{0} \times \Str{0}{b} \mid V(p)$ with
\[
V(p) = b + q + (p-1)\,d + p\,\omega m + \rho\,\bigl(h + (p-1)(h + b + q) + d\,(p-1)(p-2)/2\bigr) .
\]
The cost per turn is $V(p)/p = E'/p + \rho dp/2 + c_0'$ with $E' = (1-\rho)(b - m)$ and
$c_0' = d + \omega m + \rho\,(h + b + q) - 3\rho d/2$. For $b > m$ and $0 < \rho < 1$ it is
least at $p = \sqrt{2E'/(\rho d)}$.
\end{corollary}

\begin{proof}
The first call costs $\rho h + b + q + \omega m$, and turn $i$ of the loop costs
$\rho\,(h + b + q + id) + d + \omega m$. Expanding $V(p)$ in powers of $p$ gives the constant
term $b - m + \rho\,(m - b)$, the linear term $p\,c_0'$, and the quadratic term
$\rho d\,p^2/2$.
\end{proof}

For $p = 1$ this is the window that slides at every turn. The fixed cost of a slide is
$E' = (1-\rho)(b-m)$, the discount lost on the window, and the same square-root law applies.
In the example $E' = \text{\exStepE}$ and the best period is \exStepBest{} turns, at
\exStepPerTurn{} units per turn, less than a third of the cost of sliding at every turn.
Compaction and the stepped window are two instances of one scheme: every $p$ turns the
context is replaced by a shorter one at a fixed cost.

\subsection{The Order of the Prompt}
\label{sec:order}

The rule \textsc{Call} gives the reduced price to one term, the stored size of the first part
of the prompt. A program that puts anything new in front of its history therefore pays the
full price for the history. We give two instances.

The first is a harness that begins every prompt with a new string $t$ of one token, such as
the time or a request number:
\begin{prog}
\nm{clock}_T \defeq \kw{iter}\;T\;\kw{from}\;w_0\;\kw{as}\;i, y. \\
\quad \letin{t}{\recv{1}}\; \letin{u}{\recv{q}} & t : \Str{0}{1},\; u : \Str{0}{q} \\
\quad \letin{(\_, \_, a)}{\call{t}{y \cat u}{m}} & \text{cost } 1 + h + id + q + \omega m \\
\quad y \cat u \cat a & {} : \Str{0}{h + (i+1)d}
\end{prog}

\begin{corollary}
\label{cor:clock}
With $w_0$ of $h$ tokens,
$\vdash_0 \nm{clock}_T : \Str{0}{h + Td} \mid T\,(1 + h + q + \omega m) + d\,T(T-1)/2$.
\end{corollary}

\begin{proof}
The invariant is $\Str{0}{h + id}$. \textsc{Call} with $t : \Str{0}{1}$ and
$y \cat u : \Str{0}{h + id + q}$ costs $1 + h + id + q + \omega m$, and
$y \cat u \cat a$ has the type $\Str{0}{h + (i+1)d}$.
\end{proof}

This is the cost of the plain loop without a cache, plus $T$. The bound is attained by a
provider that grants only what it must, when the strings $t$ differ from turn to turn
(Section~\ref{sec:tests}): no prompt begins with a breakpoint. The documentation of the
provider warns of exactly this case~\cite{anthropic2026}. In the type system it is visible in
the type of the context, whose stored size stays $0$.

The second instance is the summary of Section~\ref{sec:compaction} with its instruction
before the history, $\call{\kappa}{y}{s}$ in place of $\call{y}{\kappa}{s}$. Its cost is
$g + (h + s + q + (p-1)d) + m + \omega s$, which is more by $(1-\rho)(h + s + q + (p-1)d)$. In
the example with $p = \exPbest$ this adds \exFirstExtra{} units to every block, \exFirstShare{}
percent of its cost. A person asked to summarize a text is given the request first. For a
model with a prefix cache the request belongs at the end.

\subsection{Delegation}
\label{sec:delegation}

An agent may hand a part of its work to another agent with a context of its own and receive
only a report. The history of the worker then never enters the context of the manager. We
analyze uniform hierarchies, and we assume that the work divides into tasks of any size whose
results fit a report; whether a given piece of work does is not a question of cost. All agents share a system prompt $\sigma$, a literal of $h$
tokens, and an instruction $\kappa$ of $g$ tokens that asks for a report. A task is a string
of at most $t$ tokens and a report a string of at most $r$ tokens.

A \emph{leaf} works for $k$ turns on its task $\tau$ and then reports:
\begin{prog}
\nm{leaf}_k(\tau) \defeq \letin{u}{\recv{q}} \\
\quad \letin{(\_, c, a)}{\call{\sigma}{\tau \cat u}{m}} & c : \Str{h + t + q}{0} \\
\quad \letin{y}{\iter{k-1}{c \cat a}{i, y}{\turn(y)}} & y : \Str{h + t + q + (k-1)d}{m} \\
\quad \letin{(\_, \_, z)}{\call{y}{\kappa}{r}}\; z & z : \Str{0}{r}
\end{prog}
A \emph{manager} with fan-out $J$ asks the model for a task, gives it to an agent
$\nm{sub}$ one level down, appends the report, and repeats; at the end it reports itself. Its
state is the pair of its context $y$ and the last report $v$:
\begin{prog}
\nm{man}_J(\tau) \defeq \letin{(\_, c, a)}{\call{\sigma}{\tau}{t}} & c : \Str{h+t}{0} \\
\quad \letin{v}{\nm{sub}(a)} & v : \Str{0}{r} \\
\quad \letx{(y, v)}{\kw{iter}\;J-1\;\kw{from}\;(c \cat a,\; v)\;\kw{as}\;i, (y, v).} & y : \Str{h + t + ie}{t} \\
\quad\qquad \letin{(\_, c, a)}{\call{y}{v}{t}}\; \letin{v'}{\nm{sub}(a)}\; (c \cat a,\; v')\;\kw{in} \\
\quad \letin{(\_, \_, z)}{\call{y}{v \cat \kappa}{r}}\; z & z : \Str{0}{r}
\end{prog}
Here $e = t + r$ is the growth of the manager's context per delegation: a task that the model
wrote and a report that came back.

\begin{lemma}
\label{lem:agents}
Let $\Gamma \vdash \tau : \Str{0}{t}$. Then
$\Gamma \vdash_0 \nm{leaf}_k(\tau) : \Str{0}{r} \mid S_0(k)$ with
\[
S_0(k) = (h + t + kd + g) + \omega\,(km + r) + \rho\,\bigl(k\,(h + t + q) + d\,k(k-1)/2\bigr) .
\]
If $x : \Str{0}{t} \vdash_0 \nm{sub}(x) : \Str{0}{r} \mid S$, then
$\Gamma \vdash_0 \nm{man}_J(\tau) : \Str{0}{r} \mid G(J) + J \cdot S$ with
\[
G(J) = (h + t + Je + g) + \omega\,(Jt + r) + \rho\,\bigl(J\,(h + t) + e\,J(J-1)/2\bigr) .
\]
\end{lemma}

\begin{proof}
For the leaf, the first call costs $h + t + q + \omega m$, since the literal $\sigma$ has
nothing stored. Turn $i$ of the loop costs $\rho\,(h + t + q + id) + d + \omega m$, and the
report costs $\rho\,(h + t + q + (k-1)d) + m + g + \omega r$. For the manager, the first call
costs $h + t + \omega t$. Turn $i$ of its loop costs
$\rho\,(h + t + ie) + t + r + \omega t + S$ and takes the invariant
$\Str{h + t + ie}{t} \times \Str{0}{r}$ to itself at $i + 1$. The report costs
$\rho\,(h + t + (J-1)e) + t + r + g + \omega r$. Collecting terms as in
Corollary~\ref{cor:block} gives $S_0$ and $G$.
\end{proof}

A \emph{uniform hierarchy} of depth $L$ has fan-outs $J_L, \ldots, J_1$ and leaves of $k$
turns: its agents at level $\ell \ge 1$ are managers with fan-out $J_\ell$ over agents of
level $\ell - 1$, and its agents at level $0$ are leaves. By the lemma its cost is at most
$S_L$, where $S_\ell = G(J_\ell) + J_\ell\,S_{\ell-1}$, and it performs
$k \cdot J_1 \cdots J_L$ turns of work. We call $S_L$ the \emph{cost bound} of the
hierarchy.

A part of the cost bound is the same for every organization. The coefficient of $k$ in $S_0$
is $\alpha = d + \omega m + \rho\,(h + t + q)$: a turn buys its message and its answer and
reads the prompt that every leaf begins with. We call the cost bound of a hierarchy minus
$\alpha$ times the number of its turns its \emph{overhead}. The overhead is what the organization of the
work costs, and it is where organizations differ. Let $O_L(T)$ be the least overhead of a
uniform hierarchy of depth $L$ that performs at least $T$ turns.

\begin{theorem}[Delegation]
\label{thm:delegation}
Let $\rho > 0$, $d \ge 1$, $t \ge 1$ and $r \ge 0$. For every $L \ge 0$ there are constants
$0 < a_L \le b_L$ such that for all $T \ge 1$
\[
a_L\,T^{\theta_L} \;\le\; O_L(T) \;\le\; b_L\,T^{\theta_L},
\qquad \theta_L = 1 + \frac{1}{2^{L+1} - 1} .
\]
\end{theorem}

The exponents are $2$, $4/3$, $8/7$, $16/15$ for $L = 0, 1, 2, 3$. Without delegation the
overhead is quadratic. Every level of management takes the excess of the exponent over $1$
from $1/n$ to $1/(2n+1)$.

\begin{proof}
Let $c = h + t + g + \omega r$, which is at least $1$. The overhead of a leaf of $k$ turns is
$S_0(k) - \alpha k = \rho d\,k(k-1)/2 + c$, which increases with $k$, so
$O_0(T) = \rho d\,T(T-1)/2 + c$. For $T \ge 1$ this lies between $\min(c, \rho d/4)\,T^2$ and
$(c + \rho d/2)\,T^2$, which proves the case $L = 0$.

The overhead of a hierarchy of depth $L \ge 1$ is $G(J)$ plus $J$ times the overhead of its
children, since the numbers of turns add up. The children must perform at least
$\lceil T/J \rceil$ turns each, and $O_{L-1}$ is nondecreasing, so
\[
O_L(T) = \min_{J \ge 1}\; \bigl( G(J) + J \cdot O_{L-1}(\lceil T/J \rceil) \bigr) .
\]
The function $G$ is quadratic in $J$ with the leading coefficient $\rho e/2 > 0$ and
$G(1) > 0$, so there are constants $0 < \beta \le \beta'$ with
$\beta J^2 \le G(J) \le \beta' J^2$ for all $J \ge 1$. Let $\theta = \theta_{L-1}$, and note
that $2\theta/(\theta + 1) = \theta_L$: for $\theta = 1 + 1/n$ this is
$(2n + 2)/(2n + 1)$.

\emph{Lower bound.} For every $J \ge 1$, by the induction hypothesis,
$G(J) + J\,O_{L-1}(\lceil T/J \rceil) \ge \beta J^2 + a_{L-1}\,T^\theta J^{1-\theta}$. For
positive reals $A$, $B$ and $x$ and the weights $\lambda = 2/(\theta+1)$ and
$\mu = (\theta - 1)/(\theta + 1)$, which sum to $1$, the inequality between the weighted
arithmetic and geometric means gives
\[
A\,x^{1-\theta} + B\,x^2 \;\ge\; \Bigl(\frac{A\,x^{1-\theta}}{\lambda}\Bigr)^{\lambda}
\Bigl(\frac{B\,x^2}{\mu}\Bigr)^{\mu} \;\ge\; A^\lambda B^\mu ,
\]
because the exponent of $x$ is $(1-\theta)\lambda + 2\mu = 0$ and
$\lambda^{-\lambda}\mu^{-\mu} \ge 1$. With $A = a_{L-1}T^\theta$ and $B = \beta$ this is
$a_{L-1}^{\lambda}\,\beta^{\mu}\,T^{\theta\lambda}$, and $\theta\lambda = \theta_L$.

\emph{Upper bound.} Choose $J = \lceil T^{\theta/(\theta+1)} \rceil$, so that
$J \le 2\,T^{\theta/(\theta+1)}$ and
$\lceil T/J \rceil \le T^{1/(\theta+1)} + 1 \le 2\,T^{1/(\theta+1)}$. Then
$O_L(T) \le \beta' J^2 + J\,b_{L-1}\,(2\,T^{1/(\theta+1)})^{\theta}
\le (4\beta' + 2^{\theta+1}\,b_{L-1})\,T^{2\theta/(\theta+1)}$.
\end{proof}

With a cache the cost bound of a hierarchy is not attained, because all agents begin with the
same system prompt: the worst evaluation costs $(1-\rho)h$ less for every agent after the
first (Section~\ref{sec:tests}). This replaces $h$ by $\rho h$ in the constant $c$ of the
proof and lowers $G$ by a constant, and the proof gives the same exponents for the cost of
the worst evaluation.

\paragraph{What the theorem says at practical sizes.} The exponents describe very long tasks.
For a hierarchy of depth $1$ that performs $T = Jk$ turns, the overhead per turn is
\[
\frac{\rho d}{2}\,(k - 1) \;+\; \frac{c + c_1}{k} \;+\; \frac{\rho e}{2}\Bigl(\frac{T}{k^2} - \frac{1}{k}\Bigr) + \frac{c}{T},
\qquad c_1 = e + \omega t + \rho\,(h + t) ,
\]
where $c + c_1$ is the fixed cost of one more leaf. For moderate $T$ the first two terms
decide, and the best leaf has about $\sqrt{2(c + c_1)/(\rho d)}$ turns: the square-root law
once more, with the cost of starting a leaf in the place of the cost of a summary. Only when
$T$ is so large that the manager's own rent matters does the best leaf grow, as
$(2eT/d)^{1/3}$. In the example with $t = 1{,}000$ and $r = 2{,}000$ the first formula gives
26 turns, and the script finds the best leaf at \exTabKHundred{} turns for $T = 100$, at
\exTabKThousand{} for $T = 1{,}000$, and at \exTabKTenK{} for $T = 10{,}000$. For depth $2$
the regime of the theorem lies far beyond any real task. Even with small constants
($h = 50$, $q = m = t = r = 20$, $g = 5$), the exponent fitted to the least overhead between
$T = 10^{12}$ and $T = 10^{15}$ is \exSlopeTwo{}, against $8/7 \approx 1.143$; the fixed costs
of the agents hide the limit until then. The theorem states where delegation leads in the
limit. For a given task the closed form $S_L$ is
the tool, and it can be minimized by search.

A compacting agent is a special case of this analysis with one worker after the other: the
block of Section~\ref{sec:compaction} is a leaf whose report is its summary, handed to the
next leaf. The difference between the two strategies is the manager, which keeps all reports
and pays rent on them.

\subsection{The Strategies Side by Side}

Table~\ref{tab:strategies} evaluates the closed forms of this section for the example. The
periods of compaction, of the stepped window and the size of the leaves are the best ones for
each $T$, found by search. Three things stand out. The cache is worth a factor of eight to
twelve on the plain loop. Among the strategies that bound the context, compaction is the
cheapest, delegation and the stepped window follow within a third, and the window that slides
at every turn costs more than four times as much. And for 100 turns the plain loop with a
cache is within a factor of $1.5$ of the best strategy, so a short task is not worth
compacting for money alone. The table compares bills, not results: the six strategies do not
show the model the same information.

\begin{table}[t]
\caption{Cost of $T$ turns in millions of units, for $h = 10{,}000$, $q = m = 500$,
$s = 2{,}000$, $g = 100$, $b = 20{,}000$, $t = 1{,}000$, $r = 2{,}000$, $\rho = 0.08$ and
$\omega = 4$. The plain loop exceeds a context window of 200,000 tokens after \exFeasible{}
turns; its costs beyond that are shown in parentheses.}
\label{tab:strategies}
\small
\begin{tabular}{@{}lrrr@{}}
\toprule
Strategy & $T = 100$ & $T = 1{,}000$ & $T = 10{,}000$ \\
\midrule
plain loop, no cache & \exTabBareHundred & (\exTabBareThousand) & (\exTabBareTenK) \\
plain loop & \exTabPlainHundred & (\exTabPlainThousand) & (\exTabPlainTenK) \\
window, sliding at every turn & \exTabWinHundred & \exTabWinThousand & \exTabWinTenK \\
window, sliding in steps & \exTabStepHundred & \exTabStepThousand & \exTabStepTenK \\
delegation, one level & \exTabDelHundred & \exTabDelThousand & \exTabDelTenK \\
compaction & \exTabCompHundred & \exTabCompThousand & \exTabCompTenK \\
\bottomrule
\end{tabular}
\end{table}

\section{Testing the Claims}
\label{sec:tests}

The proofs of this paper are on paper. To guard against errors in them and in the formulas,
the artifact contains a script, \texttt{cost.py}, of about 700 lines of Python without
dependencies. It has four parts.

\paragraph{An interpreter for Figure~\ref{fig:semantics}.} Strings are tuples of tokens and the
cache is a set of tuples. The interpreter has two modes. In the \emph{worst} mode every
message and every answer has its maximal length and consists of tokens that occur nowhere
else, the provider bills at the reduced price only the longest breakpoint, and a wait empties
the cache. In the \emph{random} mode lengths are drawn at random, most tokens are equal so
that prompts often begin alike, the provider bills a random number of tokens at the reduced
price above the longest breakpoint, a wait removes a random subset of the cache, and a choice
takes a random branch.

\paragraph{A checker for Figure~\ref{fig:typing} on closed programs.} A closed program has
numbers for all its indices. The checker unrolls loops, so the type that it computes for turn
$i$ is the invariant at $i$, and it joins the branches of a choice by \textsc{S-Widen}. It
returns the type, the cost, the effect, and the largest prompt.

\paragraph{The tests.} For each of \exPrograms{} random programs, four pairs of prices and
four evaluations (one in the worst mode and three random ones), the script checks the three
statements that Theorem~\ref{thm:sound} and Corollary~\ref{cor:window} make about a closed
program: the cost is within the bound, the result is in the interpretation of its type at the
final cache, and no call exceeds the largest prompt of the checker. For each of \exExact{}
random programs without choice and without $\kw{tail}$, at $\rho = 1$, it checks that the
evaluation in the worst mode costs exactly the bound (Proposition~\ref{prop:attained}). The
programs are generated from the grammar with nesting depth four and include loops whose
bodies depend on the turn number, loops over pairs, and waits. All checks pass.

\paragraph{The closed forms.} The programs of Section~\ref{sec:laws} are written as syntax
trees. For \exCases{} combinations of parameters and prices the script checks that the checker
computes exactly the closed form of the corresponding corollary and that the evaluation in
the worst mode stays within it. For the plain loop, the session, compaction, both windows and
the clock, the worst evaluation costs exactly the closed form, also for $\rho < 1$: these
bounds are attained. Two bounds are not attained with a cache. With the instruction of the
summary in front, the instruction itself is a breakpoint from the second block on, which
saves $(1-\rho)g$ per block; for an empty instruction the two orders send the same prompts,
and the rule for calls, which discounts only the first part, is loose by the whole discount.
And in a hierarchy all agents begin with the same system prompt,
so every agent after the first finds it stored: the worst evaluation costs exactly
$(1-\rho)h$ less than the bound for each of them. A program that passed the stored system
prompt down to its workers would have a type that says so. The script also checks the square-root laws against a search over whole periods, and it
computes the least overhead of hierarchies over a grid of fan-outs, for the small constants
of Section~\ref{sec:delegation}: between $T = 10^{12}$ and $T = 10^{15}$ the fitted exponents
are \exSlopeZero, \exSlopeOne{} and \exSlopeTwo{} for
depths $0$, $1$ and $2$.

A test of this kind shows only that the theorems are not false on the sampled programs. It
also shaped one statement. A first version of Theorem~\ref{thm:delegation} bounded the cost
itself, which is true but says little, because the linear part of the cost hides the exponent
up to very large $T$. The fit showed this, and the theorem now bounds the overhead. The command
\texttt{make -C artifact check} runs everything in a few seconds and compares the output,
including all numbers quoted in the text, with the committed copy.

\section{Related Work}
\label{sec:related}

\paragraph{Sizes and costs in types.} Sized types~\cite{hughes1996} and indexed
types~\cite{xi1999} put bounds on the sizes of data into types. Type systems for resource
bounds add a cost that is a function of those sizes: Crary and Weirich~\cite{crary2000}
certify running time with a virtual clock, Danielsson~\cite{danielsson2008} annotates the
types of lazy data structures with their time complexity, and Dal Lago and
Gaboardi~\cite{dallago2011} give a system with index terms that is relatively complete for
the cost of evaluation. Amortized analysis assigns potential to
data~\cite{hofmann2003,hoffmann2011} and infers polynomial bounds automatically, and
$\lambda$-amor~\cite{rajani2021} unifies several of these systems. Wegbreit~\cite{wegbreit1975}
derived closed forms for the cost of programs from recurrences fifty years ago. Our rules
stand in this tradition, and their treatment of sizes, costs and loops is standard. What is
specific is the cost model. The cost of a call is the size of the data sent, so costs are
quadratic where a count of steps would be linear, and the price of a token depends on what
was sent before.

\paragraph{State that only grows.} The stored part of a string type is a fact about the
provider's state. It stays true because that state only grows, in the sense of monotonic
state with witnessed stable predicates~\cite{ahman2018}; ``this string is a breakpoint'' is
such a predicate. Our setting adds the points where monotonicity fails. The cache forgets
during a wait, and an effect in the style of effect systems~\cite{lucassen1988} marks the
expressions after which the witnesses are void. Analyses of processor caches for bounds on
execution time~\cite{ferdinand1999} compute which memory blocks must be in the cache at a
program point; the stored size is information of this kind for a cache of prefixes.

\paragraph{Growth of loop programs.} Ben-Amram, Jones and Kristiansen~\cite{benamram2008}
decide for a language with bounded loops whether the values of variables grow polynomially.
Our index terms are polynomials, so a typable loop has polynomial cost, and the doubling loop
has no type. We do not infer invariants; their techniques may apply.

\paragraph{Programs around language models.} ReAct~\cite{yao2022} established the loop of
calls and tool results that we model. MemGPT~\cite{packer2023} manages the context like
virtual memory, with eviction and summaries, which are the strategies of
Sections~\ref{sec:compaction} and~\ref{sec:window}. Reusing the computation for a shared
prefix of prompts was proposed as Prompt Cache~\cite{gim2023}. Lumer et
al.~\cite{lumer2026} measure what the caches of three providers save on long agent tasks
under three caching strategies, and TokenPilot~\cite{xu2026} observes that pruning a context
changes its layout and invalidates the cache, and manages contexts so that prefixes stay
stable. These are the trade-offs of Sections~\ref{sec:compaction} to~\ref{sec:order}, found
by experiment; what we add is closed forms and a proof. FrugalGPT~\cite{chen2023} reduces
the cost of using models by choosing among them and adapting prompts, and
LMQL~\cite{lmql2023} shows that a programming language for prompts can save billed tokens.
There are tools that scan a code base for calls to model interfaces and estimate their price
during code review~\cite{tokentoll}. We know of no type system, and of no static analysis
with a proof of soundness, for the token cost of programs that call a model. In this series of papers, an earlier one grades such programs by their probability
of error~\cite{tpp3}; there a retry divides the error, and here it would multiply the cost.

\paragraph{Checkpoints.} The square-root law of Corollary~\ref{cor:sqrt} has the form of
Young's formula for the interval between checkpoints~\cite{young1974}, and of the older
formula of Harris for the size of a production lot, which balances the cost of setting up
against the cost of holding stock~\cite{harris1990}. In all
three a fixed cost per event competes with a cost that grows linearly between events.

\section{Limits and Open Problems}
\label{sec:limits}

\paragraph{Limits.} The bounds are statements about the semantics of
Figure~\ref{fig:semantics}, and we have not compared them with invoices. Real price lists
have features outside the model: a cache that looks back over a bounded number of blocks and
ignores short prompts, prices that depend on the length of the prompt, discounts for batches,
and tokens for reasoning that are billed and not shown. Strings are sequences of tokens,
which ignores the seams between strings. The type system bounds the worst case; an agent that
usually stops after ten turns of a possible hundred is charged for a hundred. The calculus
has no functions and no recursion, and invariants are not inferred. The laws concern money
and not the quality of the work. The proofs are not mechanized.

\paragraph{Open problems.} We list the problems that we would work on next, for whoever
continues, human or model.

\begin{enumerate}
\item \emph{Time.} With a clock in the semantics, a cache entry would expire after a lifetime
$\lambda$ and be renewed by every call. A program could then keep its cache alive during a
wait by calls with an empty answer, at $\rho P$ each. The question is a type system that
proves a call arrives in time, and the resulting law: keeping the cache alive should pay for
waits shorter than about $(1-\rho)/\rho$ lifetimes.
\item \emph{Prefixes.} The types record that a value is stored and not that one value is a
prefix of another. A program that rebuilds a stored string from its parts gets a loose bound.
Types that carry the structure of concatenations, as words over variables for segments, would
make Proposition~\ref{prop:attained} hold with a cache.
\item \emph{Inference.} The body of a loop without choice maps the sizes of its state by an
affine function with natural coefficients. The invariant is the solution of a linear
recurrence, which is polynomial unless the body duplicates its state. An algorithm that
solves it, in the manner of amortized analysis, would remove the annotations.
\item \emph{Expected cost.} With a distribution over the lengths of answers and over the turn
at which the model stops, the type should bound the expected bill. This connects to graded
types for the reliability of model calls~\cite{tpp3}: a single judgment could give the price
of a retry and what it buys.
\item \emph{Quality.} A cost for what a summary loses would turn the square-root law into a
trade between money and results. We do not know what its unit should be.
\item \emph{A real harness.} The calculus is close to the core of harnesses written in
general-purpose languages. A checker for a subset of such a language that reports the cost of
a loop as a formula, and the lines that forfeit the cache, would test whether these laws
predict bills.
\item \emph{Mechanization.} The proof of Theorem~\ref{thm:sound} is a first-order induction
with a monotone invariant, and it should take a few hundred lines in a proof assistant.
\end{enumerate}

\section{Conclusion}

The bill of an agent loop is determined by the sizes of the strings that it sends and by which
of them the provider has stored. Both are static properties of the program. A type system
with two sizes for each string, a rule for calls that returns what is now stored, and an
effect for waiting bounds the bill for every behavior of the model, and typing the common
strategies gives their costs in closed form: rent on the history, the price of a question,
the best period for compaction, the price of a sliding window, and the exponents of
delegation.

\bibliographystyle{ACM-Reference-Format}
\bibliography{refs}

\end{document}
